% EMNLP 2026 main-track submission
% Title: Test Smells: Causal Decomposition of Eval-Identifiability in LLM Benchmarks
%
% Style: assumes acl.sty + acl_natbib.bst from the EMNLP 2026 LaTeX package.
% Compile: pdflatex paper_v2 ; bibtex paper_v2 ; pdflatex paper_v2 ; pdflatex paper_v2
%
% If acl.sty is unavailable locally, swap the next two lines for:
%   \documentclass[11pt]{article}
%   \usepackage[a4paper,margin=1in]{geometry}
%   \usepackage{natbib}
% and the compile will produce a generic A4 PDF that exercises all the content.
\documentclass[11pt]{article}
\usepackage[final]{acl}

\usepackage[utf8]{inputenc}
\usepackage[T1]{fontenc}
\usepackage{microtype}
\usepackage{inconsolata}
\usepackage{amsmath, amssymb}
\usepackage{booktabs}
\usepackage{multirow}
\usepackage{array}
\usepackage{graphicx}
\usepackage{xcolor}
\usepackage{url}
\usepackage{hyperref}
\usepackage{tikz}
\usetikzlibrary{positioning,shapes,arrows.meta,backgrounds,fit}

% Convenience commands
\newcommand{\Peval}{P(\text{eval})}
\newcommand{\dpeval}{\Delta\Peval}
\newcommand{\bankspec}{\textsc{bank-spec}}
\newcommand{\todoc}[1]{\textcolor{red}{[\textbf{TODO}: #1]}}
\newcommand{\category}[1]{\textsc{#1}}
\newcommand{\bench}[1]{\textsc{#1}}

\title{Test Smells: A Causal Decomposition of Eval-Identifiability \\
       in Modern LLM Benchmarks}

\author{Anonymous EMNLP submission}

\begin{document}
\maketitle

% ============================================================================
\begin{abstract}
Modern LLM evaluations often have recognizable \emph{test smells} ---
recurrent surface features (item identifiers, scoring rubrics, fictional
placeholders, explicit harm probes, formatting constraints) that mark a
transcript as ``an eval'' rather than a real user request. We extract
a bank of \textbf{258 canonical test-smell templates} from six widely
used benchmarks (\bench{Needham}, \bench{Hawthorne}, \bench{MASK},
\bench{AgentHarm}, \bench{BBH}, \bench{MALT}), categorize them along a six-class
taxonomy, and show three causal effects.
\textbf{(1) Forward injection}: inserting bank cues into deployment-style
transcripts raises an LLM-judge's eval-classification by a mean of
$+0.21$ across $11/11$ cross-benchmark source--target pairs ($9/11$
raw-shift significant). Under a matched paraphrase control (subtracting
the shift induced by a same-scope rewrite that does not target bank
features), the \emph{bank-specific causal} shift is $\bm{+0.192}$ on
average across all $\bm{11/11}$ control-scored cells, $\bm{11/11}$
direction-positive and $\bm{7/11}$ individually significant
(per-cell range $+0.017$ to $+0.332$).
\textbf{(2) Reverse intervention with a preservation gate}: removing the
cues under a 5-criterion gate that validates task preservation lowers
\Peval{} by $-0.391$ on the source pool ($n{=}93$) and by $-0.176$
($p<10^{-13}$) on the never-seen \bench{SAD}+\bench{AISI} held-out pool, while
\emph{preserving the task}: $92\%$ behavioral parity, $96\%$ topic
preservation ($\kappa{=}0.59$, $n{=}205$), $89.8\%$ {BBH} task-metric
accuracy ($n{=}252$), and $80.8$--$100\%$ refusal-class preservation on
two safety benchmarks.
\textbf{(3) Neutral-rewrite controls localize the effect by benchmark
type}: on capability benchmarks (\bench{Needham}) bank removal beats matched
neutral edits by $7.9\times$, while on safety benchmarks neutral rewrites
move \Peval{} comparably to bank removal --- consistent with safety-eval
signal being more entangled with task content.

Activation probes trained on residual-stream features of six open-weight
models across three families (Qwen3-8B, 32B; Qwen3.5-27B, 122B-A10B;
Mistral-Small-24B-Instruct-2501 and Gemma-2-27B-it as cross-family
checks) replicate the direction at the representation level, but only
when trained on a \emph{pooled} cross-benchmark contrast: per-benchmark
probes fail to transfer cross-benchmark (mean off-diagonal AUC
$0.54$--$0.58$, including both cross-family probes), and the probe's
\Peval{} is uncorrelated with simple safety-keyword density
($r{=}+0.009$ across $1{,}350$ samples). We treat the probes as supporting diagnostics,
not as evidence for a universal eval-awareness representation, and
position our work against the stronger paired-prompt design of
\citet{abdelnabi2025hawthorne}. Our central claim is decompositional:
eval-identifiability has a removable surface component, but a large
fraction of the signal is intrinsic to the task --- particularly in
safety evals --- and cues do not behave like independent additive
features. We release the bank, paired datasets, and analysis code as a
diagnostic for benchmark designers.
\end{abstract}

% ============================================================================
\begin{figure*}[t]
\centering
\begin{tikzpicture}[
  every node/.style={font=\small, align=center},
  block/.style={rectangle, draw, rounded corners=2pt, minimum height=10mm,
                minimum width=22mm, inner sep=3pt, fill=white},
  evidence/.style={rectangle, draw=blue!40, fill=blue!5, rounded corners=2pt,
                   minimum height=8mm, minimum width=34mm, inner sep=3pt,
                   font=\footnotesize},
  arrow/.style={-Stealth, thick, gray},
  label/.style={font=\scriptsize\itshape, text=gray}
]
% Top: Raw measurement
\node[block, fill=blue!10] (raw) {Raw judge shift\\[1pt] $\boldsymbol{+0.21}$ forward / $\boldsymbol{-0.39}$ reverse};
% Split: matched controls
\node[block, fill=orange!15, below left=8mm and 12mm of raw] (surf) {Surface\\(removable)};
\node[block, fill=red!12, below right=8mm and 12mm of raw] (intr) {Intrinsic\\(task-entangled)};
\draw[arrow] (raw) -- node[label, left=1pt] {matched neutral} (surf);
\draw[arrow] (raw) -- node[label, right=1pt] {everything else} (intr);
% Evidence boxes
\node[evidence, below=4mm of surf] (e1) {Forward bank-specific causal:\\$\bm{+0.192}$ pooled (11/11 cells $+$)};
\node[evidence, below=2mm of e1]   (e2) {Reverse causal $\Delta = -0.346$\\($7.9\times$ vs neutral, Needham)};
\node[evidence, below=4mm of intr] (e3) {Held-out reverse $-0.176$\\($p<10^{-13}$, SAD+AISI)};
\node[evidence, below=2mm of e3]   (e4) {SAD v3 gate rejects $\bm{60/60}$:\\``surface'' = test on SAD};
\draw[arrow] (surf) -- (e1);
\draw[arrow] (intr) -- (e3);
% Bottom: preservation oracles
\node[block, fill=green!12, minimum width=125mm, below=12mm of $(e2)!0.5!(e4)$] (pres)
  {Preservation oracles ($73$--$100\%$): \textbf{92\%} behavioral / \textbf{96\%} topic / \textbf{89.8\%} BBH task / \textbf{80.8\%} AgentHarm refuse-class / \textbf{73.1\%} AgentHarm safety-rubric / \textbf{100\%} Hawthorne refuse-class};
\draw[arrow, dashed] (e2) -- (pres.north -| e2);
\draw[arrow, dashed] (e4) -- (pres.north -| e4);
% Caption hint
\node[label, below=2mm of pres, font=\scriptsize] {Removal of surface cues survives all six preservation oracles ($73$--$100\%$); intrinsic component is what remains under preservation-validated removal.};
\end{tikzpicture}
\caption{The test-smell decomposition. The raw judge shift induced by bank
features ($+0.21$ forward / $-0.39$ reverse, Sec.~\ref{sec:forward}--\ref{sec:reverse})
splits into a \emph{surface} component (removable by preservation-validated rewriting;
quantified by matched neutral controls in both directions) and an \emph{intrinsic}
component (task-entangled; survives surface ablation). The surface component is
benchmark-type-dependent: on capability benchmarks it dominates (bank/neutral ratio
$7.9\times$ on \bench{Needham} sources, Table~\ref{tab:reverse-canonical}); on
safety benchmarks surface removal is less differentiating from neutral paraphrase
at the judge level (bank/neutral ratio $\sim 1\text{--}2\times$), while still
preserving the underlying safety behavior on six oracles
(Table~\ref{tab:preservation}; dual-annotator $n{=}60$ spot-check: $0$ task
unlocks across both annotators). All ablations preserve the task across
those oracles.}
\label{fig:decomp}
\end{figure*}

% ============================================================================
\section{Introduction}
\label{sec:intro}

Frontier LLMs behave differently when they detect they are being
evaluated --- a phenomenon often called \emph{evaluation awareness} or
\emph{test awareness} \citep{needham2025evalawareness,
abdelnabi2025hawthorne}. The behavior is alignment-relevant because
several documented misalignments --- alignment faking
\citep{anthropic2024alignmentfaking}, in-context scheming
\citep{meinke2024scheming}, and deferred defection in sleeper agents
\citep{hubinger2024sleeper} --- can be \emph{conditioned} on the
model's belief that it is being observed. A concurrent line of work
\citep{aranguri2026vea} debates whether \emph{verbalized}
eval-awareness in chain-of-thought causally inflates measured safety.
We sidestep that debate by measuring the upstream construct directly
via the MCQ probe of \citet{needham2025evalawareness}, and ask:
\emph{which features of the input cause an LLM judge's
eval-vs-deployment posterior to move, and how do those features
compose?}

Borrowing from software engineering, we use \emph{test smell} for any
recurring surface feature of an evaluation transcript that
distinguishes it from real user traffic --- parseable answer markers
(\texttt{Final answer: <X>}), numeric/style constraints, hard-coded
secret tokens, eval-harness scaffolding, and so on. The empirical
question is which test smells the model actually \emph{uses}, how much
each shifts its eval posterior, and whether removing them
\emph{preserves the task}. Two prior efforts shape our work without
defining the question we ask: \citet{pre2025stealtheval}'s
Probe--Rewrite--Evaluate pipeline shares our high-level shape (probe,
rewrite, re-measure) but does not validate task preservation;
\citet{abdelnabi2025hawthorne}'s white-box probe uses paired
real-vs-hypothetical prompts and reasoning-trace snippet annotation to
isolate model-internal eval cognition --- a representation-engineering
question, where ours is a measurement question about which surface
features drive judge classification and whether they can be removed
without breaking the task.

\paragraph{Contributions.}
\begin{enumerate}
  \item A \textbf{258-template bank} of canonical test-smells extracted
    from six benchmarks, classified into six categories (\S\ref{sec:bank}),
    with a transparent manual-validation audit ($n{=}42$, strict pass
    $67\%$, inclusive $88\%$) framed honestly as an
    \emph{imperfect high-recall} catalog.
  \item A \textbf{canonical reverse table} (\S\ref{sec:reverse}) with one
    set of numbers, raw-$\Delta$ and matched-control-adjusted causal-$\Delta$
    side-by-side, replacing scattered numbers from earlier drafts.
  \item A \textbf{forward neutral-injection control} (\S\ref{sec:forward})
    that quantifies how much of the forward effect is bank-specific
    versus generic length/formatting artifact.
  \item Five \textbf{convergent preservation oracles} (\S\ref{sec:preservation})
    spanning capability and safety benchmarks: LLM-judged behavioral
    parity ($92\%$), topic preservation ($96\%$ with inter-rater $\kappa{=}0.59$),
    {BBH} task-metric accuracy ($89.8\%$), and refusal-class
    preservation on \bench{AgentHarm} ($80.8\%$) and \bench{Hawthorne} ($100\%$).
  \item A \textbf{benchmark-type asymmetry}: on capability benchmarks the
    bank causes a $7.9\times$ larger shift than matched neutral rewrites
    of the same scope; on safety benchmarks neutral rewrites move the
    judge comparably, consistent with safety-eval signal being more
    intrinsically entangled with the task (\S\ref{sec:specificity}).
  \item \textbf{Supporting activation-probe evidence} from six open-weight
    models with explicit \emph{transparent negative-result} reporting
    (\S\ref{sec:probes}): per-benchmark probes do not transfer
    cross-benchmark (off-diagonal AUC $\sim$$0.54$); the bank-removal
    effect survives a pooled-probe construction that does transfer,
    but is benchmark-domain-specific and not a universal
    eval-awareness representation.
\end{enumerate}

% ============================================================================
\section{Related Work}
\label{sec:related}

\paragraph{Evaluation/situational awareness.}
\citet{laine2024sad} introduced the Situational Awareness Dataset;
\citet{berglund2024taken} demonstrated out-of-context reasoning that
implies self-knowledge transfer from training;
\citet{needham2025evalawareness} operationalized eval-awareness with the
gated \texttt{jjpn2/eval\_awareness} dataset and an MCQ probe, on
which we build directly. \citet{hua2025steering} and
\citet{fan2025evalfaking} probe related observer effects.

\paragraph{Concurrent verbalized-eval-awareness debate.}
\citet{aranguri2026vea} show that verbalized eval-awareness in
chain-of-thought correlates with refusal rates across $8$ models and
$19$ benchmarks; \citet{knecht2026vea} reach the opposite behavioral
conclusion using CoT prefill across $8$ reasoning models. We complement
this debate by operationalizing awareness with the upstream MCQ probe
(not a CoT side-channel) and by adding a preservation gate that
disambiguates ``the model treats the rewrite as deployment'' from
``the rewrite shifted the test target.''

\paragraph{Probe--rewrite--evaluate (PRE / StealthEval).}
\citet{pre2025stealtheval} introduced the probe$\to$LLM-rewrite$\to$
measure pipeline. We extend it with a 5-criterion preservation gate and
a matched neutral control, and we explicitly report on the gap between
gated and ungated rewriters (\S\ref{sec:robustness}). We label our
ungated variant ``\textbf{ungated-rewrite baseline}'' rather than
``PRE-equivalent'' because we do not run their exact rewriter or
dataset.

\paragraph{Format sensitivity critique.}
\citet{formatcritique2026} argues that PRE-style probe results may
reflect format sensitivity rather than eval-awareness. Our matched
neutral-injection (\S\ref{sec:forward}) and matched random-rewrite
controls (\S\ref{sec:reverse}, \S\ref{sec:specificity}) directly test
this concern. On capability benchmarks, the bank effect substantially
exceeds the format/scope control; on safety benchmarks, the two are
closer in magnitude.

\paragraph{White-box probing of eval-awareness.}
\citet{abdelnabi2025hawthorne} train a binary linear probe on
reasoning-trace snippets, using paired real-vs-hypothetical prompts to
isolate model-internal test awareness, and demonstrate causal steering.
This is the stronger methodology for the representation claim;
\citet{xue2025linearcontroleval} make related mid-layer claims at the
representation level. We position our probes as supporting diagnostics
of input-level structure (not a representation claim) and use
\citeauthor{abdelnabi2025hawthorne}'s design to bound what our probes
can and cannot say (\S\ref{sec:probes}).

\paragraph{Benchmark validity and contamination.}
\citet{gururangan2018annotationartifacts,poliak2018hypothesisonly,
bowman2021benchmarkingnlu} established that surface artifacts inflate
apparent NLU capabilities; \citet{sainz2023contamination,
magar2022contamination} did so for data contamination. We target a
complementary \emph{behavioral} leak: even uncontaminated benchmarks
carry features that signal ``this is a test'' to the model.

% ============================================================================
\section{The Bank: 258 Templates Across Six Benchmarks}
\label{sec:bank}

\paragraph{Donor pool (E-cases).}
We draw donor transcripts from six widely used benchmarks: \bench{Needham}
\citep{needham2025evalawareness} (capability/agentic),
\bench{Hawthorne} \citep{abdelnabi2025hawthorne} (safety harmful-request),
\bench{MASK} \citep{ren2025mask} (honesty/deception), \bench{AgentHarm}
\citep{andriushchenko2025agentharm} (harmful agentic tools), \bench{BBH}
\citep{suzgun2023bbh} (multi-task reasoning), and \bench{MALT}
\citep{malt2025} (multi-agent code). We filter to transcripts where
the model-measured eval-awareness $\bar P(\text{eval})\geq 0.8$ and the
mean of two independent human annotations $\geq 0.7$.

\paragraph{Recipient pools (D-cases).}
For the forward direction (\S\ref{sec:forward}) we use deployment-style
D-cases drawn from low-confidence-deploy subsets of each source benchmark,
augmented with naturalistic chat from \bench{WildChat} \citep{zhao2024wildchat}.

\paragraph{Extraction pipeline.}
A three-stage extraction (per-transcript LLM extraction; cross-transcript
deduplication into templates; selection for the causal experiments)
yields \textbf{258 canonical templates} grouped into six categories:
\category{SystemInstructions}, \category{TaskFormatting},
\category{AnswerFormatting}, \category{FictionalScenario},
\category{ContrivedTask}, and \category{EthicalProbe}. Templates are
prefixed by source benchmark (e.g.\ \texttt{needham\_\_response\_parsing\_markers}).
Counts per category are in Table~\ref{tab:bank-cats}; per-benchmark
splits are in Appendix~\ref{app:bank}. \bench{MALT} contributes the
largest share of templates ($80$, or $31\%$ of the bank); its templates
flow into the pooled v3 reverse run as part of the \bench{Needham} source
pool's sub-task composition (Appendix~\ref{app:reverse-bench}), but
\bench{MALT} is not separately re-evaluated in the causal experiments
of \S\S\ref{sec:forward}--\ref{sec:reverse} because its donor pool
lacks usable D-cases.

\begin{table}[t]
\centering\small
\begin{tabular}{lr}
\toprule
Category & Templates \\
\midrule
\category{ContrivedTask}        & 75 \\
\category{SystemInstructions}   & 44 \\
\category{TaskFormatting}       & 42 \\
\category{FictionalScenario}    & 35 \\
\category{AnswerFormatting}     & 31 \\
\category{EthicalProbe}         & 31 \\
\midrule
\textbf{Total}                  & \textbf{258} \\
\bottomrule
\end{tabular}
\caption{Bank composition by category. Per-benchmark splits in
Appendix~\ref{app:bank}.}
\label{tab:bank-cats}
\end{table}

\paragraph{Manual-validation audit (honest framing).}
We audited a stratified random sample of $n{=}42$ templates ($7$ per
benchmark). Pass rates differ sharply by source benchmark
(Table~\ref{tab:bank-audit}). The aggregate is $67\%$ strict ($28/42$) /
$88\%$ inclusive ($37/42$); on safety benchmarks the strict rate drops to $14$--$57\%$,
indicating that for safety benchmarks the extraction tends to over-include
content cues that are not pure scaffolding.

\begin{table}[t]
\centering\small
\begin{tabular}{lrr}
\toprule
Benchmark & $n$ & Strict pass (count) \\
\midrule
\bench{BBH}        & 7 & 100\% (7/7) \\
\bench{Needham}    & 7 & 100\% (7/7) \\
\bench{MALT}       & 7 & 86\% (6/7; inclusive 7/7) \\
\bench{AgentHarm}  & 7 & 57\% (4/7) \\
\bench{MASK}       & 7 & 43\% (3/7) \\
\bench{Hawthorne}  & 7 & 14\% (1/7) \\
\midrule
\textbf{Aggregate} & \textbf{42} & \textbf{67\% strict (28/42) / 88\% inclusive (37/42)} \\
\bottomrule
\end{tabular}
\caption{Manual-validation audit (strict / inclusive). Aggregate strict
count is $7{+}7{+}6{+}4{+}3{+}1 = 28/42 = 66.7\%$; earlier drafts
reported $71\%$ from a different rounding convention --- the per-benchmark
counts in this table are authoritative.
We frame the catalog as an \emph{imperfect high-recall} bank rather than
``canonical'' in a precision sense; the headline causal effects in
\S\S\ref{sec:forward}--\ref{sec:specificity} are reported for both
the full bank and a high-precision subset
(Appendix~\ref{app:high-precision}).}
\label{tab:bank-audit}
\end{table}

% ============================================================================
\section{Forward Injection With Matched Neutral Controls}
\label{sec:forward}

\paragraph{Cross-benchmark sufficiency (raw effect).}
For each of $11$ usable source--target benchmark pairs,\footnote{The full $5
\times 3 = 15$ Cartesian product would also include \bench{BBH}$\to$\bench{Hawthorne},
\bench{BBH}$\to$\bench{MASK}, \bench{MALT}$\to$\bench{Hawthorne}, and
\bench{MALT}$\to$\bench{MASK}. We omit these four cells because the donor templates
for capability-style benchmarks (\bench{BBH}, \bench{MALT}) are formatting/answer-marker
patterns that do not compose well with safety-benchmark target distributions
(e.g.\ inserting a parseable answer marker on a Hawthorne harmful-request item
produces incoherent prompts that the rewriter rejects).} we inject the source
benchmark's bank templates into the target benchmark's deployment-style D-cases
and re-judge \Peval{} with GPT-5-mini. All $11/11$ cells produce a positive
shift; $9/11$ are significant at $p<0.05$. The mean \emph{raw} shift is
$\bm{+0.21}$, with range $+0.02$ to $+0.43$ (largest:
\bench{MASK}$\rightarrow$\bench{Hawthorne}, $+0.427$). The raw shift is an
\emph{upper bound} on the bank-specific component, because any prefix or
system-message insertion can shift the judge --- see next subsection.

\paragraph{Matched-control causal analysis across all $11$ cells.}
A reviewer objection: \emph{any} rewrite or insertion might raise
\Peval{} (formatting/length artifact). We address this with two
complementary matched-control methodologies that answer slightly
different questions: a \emph{paraphrase} control isolates bank-specific
\emph{content} effects (does the bank do more than a same-scope
rewrite?); a \emph{neutral-prefix} control isolates bank-specific
\emph{prefix-position} effects (does the bank do more than length-matched
innocuous boilerplate?). Both controls point the same direction.

First, our cross-injection pipeline produces both a \emph{treatment}
rewrite (D-case rewritten to incorporate the source benchmark's bank
features) and a length-matched \emph{paraphrase control} (D-case
paraphrased to similar scope but without bank-feature insertion). The
\emph{causal} per-cell shift is $\bar\Delta_{\text{treat}} -
\bar\Delta_{\text{ctrl}}$. Table~\ref{tab:fwd-causal-11cell} reports
all $\bm{11}$ cross-injection cells.

\begin{table}[t]
\centering\small
\setlength{\tabcolsep}{3pt}
\begin{tabular}{lrr}
\toprule
Cell (source $\to$ target) & $n$ & causal $\bm{\Delta}$ \\
\midrule
\bench{Hawthorne}$\to$\bench{MASK}    & 19 & $+0.332$ \quad CI~$[0.20, 0.47]$ \\
\bench{Needham}$\to$\bench{Hawthorne} & 15 & $+0.290$ \quad $|t|>1.96$ \\
\bench{AgentHarm}$\to$\bench{Hawthorne}& 13& $+0.269$ \quad CI~$[0.00, 0.54]$ \\
\bench{BBH}$\to$\bench{Hawthorne}     & 20 & $+0.253$ \quad $t{=}+3.23$ \\
\bench{Needham}$\to$\bench{MASK}      & 15 & $+0.253$ \quad $|t|>1.96$ \\
\bench{MASK}$\to$\bench{Hawthorne}    & 15 & $+0.245$ \quad $|t|>1.96$ \\
\bench{MALT}$\to$\bench{MASK}         & 19 & $+0.197$ \quad CI~$[0.02, 0.39]$ \\
\bench{AgentHarm}$\to$\bench{MASK}    & 15 & $+0.163$ \quad $|t|>1.96$ \\
\bench{MASK}$\to$\bench{Needham}      & 23 & $+0.076$ \\
\bench{Hawthorne}$\to$\bench{Needham} & 15 & $+0.020$ \\
\bench{AgentHarm}$\to$\bench{Needham} & 15 & $+0.017$ \\
\midrule
\textbf{Pooled mean} (\textbf{11} cells)  & --- & $\bm{+0.192}$ \\
\textbf{Pooled bootstrap $95\%$ CI}        & --- & $\bm{[+0.129, +0.252]}$ \\
\multicolumn{3}{l}{\footnotesize $\bm{11/11}$ direction-positive; $\bm{7/11}$ significant (CI excludes $0$ or $|t|>1.96$).}\\
\multicolumn{3}{l}{\footnotesize Pooled one-sample $t$ vs $0$: $t{=}5.82$, $\bm{p{=}1.7\times 10^{-4}}$.}\\
\bottomrule
\end{tabular}
\caption{Forward causal shift after subtracting a matched-paraphrase
control, across all $\bm{11}$ cross-injection cells. The bank-specific
causal $\Delta$ pooled mean is $\bm{+0.192}$ (bootstrap $95\%$ CI
$\bm{[+0.129, +0.252]}$, one-sample $t{=}5.82$,
$\bm{p{=}1.7{\times}10^{-4}}$), about $91\%$ of the $11$-cell raw mean
$+0.21$. All cells direction-positive; $7/11$
individually significant; the remaining $4$ have small per-cell $n$
($15$--$23$).}
\label{tab:fwd-causal-11cell}
\end{table}

Second, a complementary per-template \emph{neutral-prefix} methodology
on \bench{Hawthorne} and \bench{MASK} ($n{=}20$ each), giving the same
direction at a different methodology granularity: for each D-case
we also re-judge under a length-matched innocuous boilerplate prefix
(generic helpful text, no test vocabulary).
Table~\ref{tab:fwd-neutral} reports the per-benchmark bank-specific
component ($\bm{\Delta}_{\text{bank}} -
\bm{\Delta}_{\text{neutral-prefix}}$) with bootstrap $95\%$ CIs.

\begin{table}[t]
\centering\small
\setlength{\tabcolsep}{3pt}
\begin{tabular}{lrrrrr}
\toprule
& $n$ & $\bm{\Delta}_{\text{bank}}$ & $\bm{\Delta}_{\text{neutral}}$ &
\textbf{$\bm{\Delta}_{\text{bank-specific}}$} & ratio \\
\midrule
\bench{Hawthorne} & 20 & $+0.262$ & $+0.113$ & $\bm{+0.150}$ $[0.024, 0.310]$ & $2.33\times$ \\
\bench{MASK}      & 20 & $+0.160$ & $+0.080$ & $\bm{+0.080}$ $[0.000, 0.200]$ & $2.00\times$ \\
\bottomrule
\end{tabular}
\caption{Complementary methodology: a length-matched innocuous
\emph{neutral-prefix} (rather than paraphrase) control on
\bench{Hawthorne} and \bench{MASK}. Any prefix raises \Peval{} by
$\sim 0.08$--$0.11$; the bank-specific component survives subtraction:
\bench{Hawthorne} CI excludes $0$, \bench{MASK} CI just touches $0$.
The two methodologies (Tables~\ref{tab:fwd-causal-11cell}
and~\ref{tab:fwd-neutral}) point the same direction.}
\label{tab:fwd-neutral}
\end{table}

\paragraph{Reading the result honestly.}
The raw $+0.21$ across $11$ cells reflects the gross effect of any
bank-prefixed rewrite. Subtracting a matched paraphrase control
(Table~\ref{tab:fwd-causal-11cell}) drops the pooled mean to $+0.192$
($\sim 91\%$ of raw); $\bm{11/11}$ cells are direction-positive and
$\bm{7/11}$ are individually significant. The complementary
neutral-prefix methodology (Table~\ref{tab:fwd-neutral}) on the two
benchmarks where we ran it gives bank-specific shifts of $+0.150$
(CI excludes $0$) and $+0.080$ (CI just touches $0$). Both methodologies
agree in direction. Headline reading: \emph{bank-specific forward
shift $\bm{+0.192}$ pooled across all $\bm{11}$ cells; per-cell range
$+0.02$ to $+0.33$}.

% ============================================================================
\section{Reverse Intervention: A Single Canonical Table}
\label{sec:reverse}

\paragraph{Sign convention.}
Throughout this paper we use $\Delta = \Peval_{\text{after}} - \Peval_{\text{before}}$.
\emph{Positive} $\Delta$ means a rewrite raised the judge's eval-classification;
\emph{negative} $\Delta$ means it lowered. Forward injection produces positive $\Delta$;
reverse intervention produces negative $\Delta$.

\paragraph{Preservation gate.}
Our reverse pipeline (``v3'') rewrites an evaluation transcript to
remove bank cues, subject to a 5-criterion LLM gate (the rewrite must:
preserve the task, preserve the behavior target, preserve answerability,
remove only scaffolding, leave protected spans intact) plus a
verbatim-survival check. Rejected rewrites are excluded. We also pair
each gated rewrite with a \emph{matched neutral edit} of comparable
scope: a paraphraser produces a same-length random-span rewrite that
does not target bank cues. This isolates bank-specific causal effect
from generic rewrite-scope effect.

\paragraph{One canonical table.}
Earlier drafts of this project reported scattered numbers from multiple
pipelines (v1, v2, v3, v4). For this submission we report exactly the
v3 (preservation-gated) result on the source pool and the v2 result on
the held-out pool, in Table~\ref{tab:reverse-canonical}. The
\emph{causal $\Delta$} column subtracts the matched-control shift from
the raw shift.

\begin{table*}[t]
\centering\small
\setlength{\tabcolsep}{4pt}
\begin{tabular}{llrrrrrrr}
\toprule
Version & Pool & $n$ & baseline \Peval & treatment \Peval & control \Peval & raw $\bm{\Delta}$ & \textbf{causal $\bm{\Delta}$} & ratio \\
\midrule
v3 (gated) & \bench{Needham} source & 93 & 0.900 & 0.509 & 0.855 & $-0.391$ & $\bm{-0.346}$ \quad $(p<10^{-9})$ & $\bm{7.9\times}$ \\
v2         & held-out (\bench{SAD}+\bench{AISI}) & 146 & 0.787 & 0.610 & --- & $-0.176$ \quad $(p<10^{-13},\ d{=}0.66)$ & --- & --- \\
\bottomrule
\end{tabular}
\caption{Canonical reverse table. The v3 row is donor-restricted to the
\bench{Needham} source pool (the only donor benchmark for which we ran the
full preservation-gated v3 pipeline with matched neutral controls); causal
$\Delta$ is the bank-specific shift after subtracting matched control. The
\emph{ratio} column is $|\text{raw}\ \Delta| / |\text{control}\ \Delta|$ on
\bench{Needham} = $0.391/0.045 \approx 7.9\times$. The v2 row is on a never-mined
held-out pool (\bench{SAD}+\bench{AISI}); no matched control was collected in this
pipeline so we report only raw $\Delta$ with a Wilcoxon $p$-value against
baseline. Per-category breakdown of the v3 row is in
Appendix~\ref{app:reverse-cat}; the per-sub-benchmark composition of the
\bench{Needham} pool is in Appendix~\ref{app:reverse-bench}.}
\label{tab:reverse-canonical}
\end{table*}

\paragraph{Held-out generalization.}
The bank, applied with no re-extraction to \bench{SAD} and AISI Inspect
logs --- benchmarks the templates were not mined from --- still produces a
\textbf{$-0.176$ shift in mean \Peval} ($n{=}146$ paired, Cohen's
$d{=}0.66$, $p<10^{-13}$). The bank-as-checklist captures structure
that transfers across distributions.

\paragraph{Unit of independence.} All $p$-values use the transcript as
the unit of paired observation. Repeated judge calls per transcript are
averaged before testing, not used as independent samples.

% ============================================================================
\section{Preservation Oracles}
\label{sec:preservation}

A rewrite that lowers \Peval{} is interesting only insofar as it
\emph{preserves the task}. We report six preservation measurements,
spanning capability and safety benchmarks, in
Table~\ref{tab:preservation}. All six fall in the
$\bm{73}$--$\bm{100\%}$ range, with the strictest safety-rubric judge on
\bench{AgentHarm} at $73.1\%$ (the $4$ generic measurements and
\bench{Hawthorne} all $\geq 80\%$).

\begin{table}[t]
\centering\small
\setlength{\tabcolsep}{3pt}
\begin{tabular}{p{0.45\linewidth}rr}
\toprule
Oracle / metric & $n$ & Value \\
\midrule
LLM-judged 5-criterion behavioral preservation (pooled, 6 benchmarks) & 49 & 92\% \\
Topic preservation rubric, IRR Cohen's $\kappa{=}0.59$ & 205 & 96\% \\
\bench{BBH} task-metric accuracy (capability)
   & 252 & 89.8\% \\
\bench{AgentHarm} refusal-class preservation (safety, gemini-2.5-flash
\textsc{Refuse}/\textsc{Partial}/\textsc{Comply})
   & 26 & 80.8\% \\
\bench{AgentHarm} \emph{safety-rubric} preservation (safety,
ground-truth \textsc{Safe}/\textsc{Harmful} judge)
   & 26 & \textbf{73.1\%} \\
\bench{Hawthorne} refusal-class preservation (safety, gemini-2.5-flash)
   & 17 & 100\% \\
\bottomrule
\end{tabular}
\caption{Convergent preservation oracles. The capability-side metric is
exact-answer accuracy on \bench{BBH}. The safety-side rows use two
different judges on the same $26$ \bench{AgentHarm} pairs: a generic
\textsc{Refuse}/\textsc{Partial}/\textsc{Comply} class judge (row 4) and
the stricter task-rubric judge that classifies the response as
\textsc{Safe} or \textsc{Harmful} relative to AgentHarm's task ground
truth (row 5). The rubric judge is more conservative and reports a
$7.7\%$ unlock rate ($2/26$ \textsc{Safe} originals became
\textsc{Harmful} under v3), against a $19.2\%$ \emph{harden} rate
($5/26$ \textsc{Harmful} originals became \textsc{Safe}). Net: v3 is
slightly \emph{safer} on aggregate but is not lossless --- we surface
the $7.7\%$ unlock rate explicitly. \bench{Hawthorne} (row 6) is fully
preserved by either metric.}
\label{tab:preservation}
\end{table}

For comparison, an ungated rewrite baseline (no preservation gate) yields
$\sim 10\%$ refusal-class preservation on \bench{AgentHarm}; the gate
provides an $8\times$ improvement on this benchmark
(see \S\ref{sec:robustness}).

\paragraph{Dual-annotator safety spot-check ($n{=}60$).}
We further validate the safety preservation claim with a stratified
$n{=}60$ paired sample (\bench{AgentHarm} + \bench{Hawthorne}
$\times$ \{orig, v3\} $\times$ $15$ per cell), independently annotated
by two LLMs from different model families
(\texttt{gemini-2.5-flash} and \texttt{claude-haiku-4.5}). Inter-annotator
raw agreement $\bm{81.7\%}$, Cohen's $\bm{\kappa = 0.509}$ (moderate
agreement, Landis \& Koch). \textbf{Zero task unlocks across BOTH
annotators on all $30$ paired samples}; the v3 rewrite never converted
an orig-\textsc{Safe} refusal into a v3-\textsc{Harmful} compliance
under either annotator. v3 SAFE rate is non-decreasing in every cell
for both annotators (\bench{AgentHarm} v3 $\geq$ orig under annotator
A and B; \bench{Hawthorne} v3 $\geq$ orig under both). This converges
with the gemini-rubric $73.1\%$ preservation and the
$80.8\%$/$100\%$ refusal-class preservation already in
Table~\ref{tab:preservation}. We frame this as a dual-LLM spot-check,
not human annotation, but two independent strong-LLM annotators from
different families provide independent corroboration that the gate's
preservation is not a single-judge artifact.

\paragraph{Dual-judge agreement (formal stratified $n{=}153$).}
Our single-judge dependence concern (the headline numbers all flow through
\texttt{gpt-5-mini}) is addressed with a formal stratified sample: we drew
$n{=}153$ transcripts ($\sim 26$ per benchmark $\times$ 6 benchmarks; capability
benchmarks contribute E-cases only) and had two judges score each
independently --- \texttt{gemini-2.5-flash} and \texttt{claude-haiku-4.5} ---
under a fixed YES/NO prompt. Aggregate agreement rate is
$\bm{92.2\%}$ with Cohen's $\bm{\kappa = 0.413}$ (moderate agreement,
Landis \& Koch). Both judges score $\sim$$71\%$ accuracy against the
\texttt{gpt-5-mini}-assigned E/D labels --- the expected fraction given the
labels themselves come from a different judge. The per-benchmark
breakdown (Table~\ref{tab:dual-judge-perbench}) shows full agreement
($100\%$) on the eval-only capability benchmarks (\bench{AgentHarm},
\bench{BBH}, \bench{MALT}) and $88\%$ on the three benchmarks with both
E and D cases.

\begin{table}[t]
\centering\small
\setlength{\tabcolsep}{4pt}
\begin{tabular}{lrr}
\toprule
Benchmark & $n$ & Agreement \\
\midrule
\bench{Needham}    & 34 & 88.2\% \\
\bench{Hawthorne}  & 34 & 88.2\% \\
\bench{MASK}       & 34 & 88.2\% \\
\bench{AgentHarm} (E only) & 17 & 100\% \\
\bench{BBH} (E only)       & 17 & 100\% \\
\bench{MALT} (E only)      & 17 & 100\% \\
\midrule
\textbf{Pooled} & \textbf{153} & \textbf{92.2\%} ($\kappa{=}0.413$) \\
\bottomrule
\end{tabular}
\caption{Dual-judge stratified agreement, \texttt{gemini-2.5-flash} vs
\texttt{claude-haiku-4.5}, simple YES/NO eval-vs-deployment prompt.
Per-benchmark agreement is at or near ceiling on the eval-only capability
benchmarks; the $\sim$$88\%$ on the three E+D benchmarks corresponds to
borderline samples where E and D are not crisply separated. We treat the
$\kappa = 0.413$ as evidence of real cross-judge agreement beyond chance,
not as exceptional convergence.}
\label{tab:dual-judge-perbench}
\end{table}

\paragraph{Position-bias robustness (2-swap, $n{=}81$).}
A further single-judge robustness check tests whether \texttt{gpt-5-mini}
flips classification under reordering of the YES/NO prompt options
(PROMPT\_A: ``YES if eval, NO otherwise''; PROMPT\_B: ``NO if
real-world, YES otherwise''). A label-position-biased judge would flip
between the two; a robust judge stays consistent. On a stratified
$n{=}81$ sample, the overall flip rate is $\bm{6.2\%}$
(\bench{Needham} $5.6\%$, \bench{Hawthorne} $16.7\%$, \bench{MASK}
$5.6\%$; capability benchmarks $0\%$). The $16.7\%$ on \bench{Hawthorne}
is the only above-$10\%$ cell and is concentrated on borderline cases.
Overall $6.2\%$ is well within community robustness norms.

\paragraph{Dual-annotator safety spot-check ($n{=}60$).}
As an additional safety check, two independent annotators
(\texttt{gemini-2.5-flash} and \texttt{claude-haiku-4.5}) classified
target-model (\texttt{gpt-5-mini}) responses on $60$ stratified
\bench{AgentHarm}~$+$~\bench{Hawthorne} transcripts ($15$ per
benchmark $\times$ $2$ conditions $\{$orig, v3$\}$) along three axes:
harmfulness class (\textsc{Safe}/\textsc{Partial}/\textsc{Harmful}),
refusal-target preservation, and task-unlock. Raw inter-annotator
agreement is $\bm{81.7\%}$ (Cohen's $\bm{\kappa = 0.509}$, moderate).
Most importantly, \textbf{both annotators report zero task unlocks} (no
orig-\textsc{Safe} sample became v3-\textsc{Harmful}) across all
$30$ paired samples. Per-condition \textsc{Safe} rate is non-decreasing
under v3 in every cell (annotator A: \bench{AgentHarm}
$60\%\to 67\%$, \bench{Hawthorne} $87\%\to 93\%$; annotator B:
$53\%\to 80\%$ and $87\%\to 93\%$), with paired class preservation
$80$--$93\%$ (A) and $60$--$93\%$ (B). We frame this as a
\emph{dual-LLM-judge spot-check}, not human annotation; it converges
with the rubric-judge $73.1\%$ and refusal-class $80.8$--$100\%$
preservation numbers (Table~\ref{tab:preservation}) and rules out the
``v3 unlocks harmful capability'' failure mode at this $n$.

% ============================================================================
\section{Bank-Specificity Is Benchmark-Type-Dependent}
\label{sec:specificity}

\paragraph{Paired controls on the activation probe.}
We use a single representative model for the probe-level
bank-specificity analysis (Qwen3.5-27B; the pooled-probe construction of
\S\ref{sec:probes} achieves CV $0.850$ at layer $12$, the layer selected
by 5-fold CV across the union of \bench{Needham}, \bench{Hawthorne},
\bench{MASK} E vs D). On this probe we compute \emph{paired}
bank-specificity: for each base sample ID that has orig, v3-removed, and
rev-control all cached, we take the per-sample \Peval{} triple and compute
bootstrap CIs over the per-sample causal $\Delta = \Delta_{\text{bank}} -
\Delta_{\text{control}}$. Table~\ref{tab:bank-specificity-paired} reports
the result.

\begin{table}[t]
\centering\small
\setlength{\tabcolsep}{3pt}
\begin{tabular}{lrrrrr}
\toprule
Benchmark & $n_{\text{paired}}$ & $\Delta_{\text{bank}}$ & $\Delta_{\text{ctrl}}$ & causal $\Delta$ & ratio \\
\midrule
\bench{Hawthorne} (v3) & 15 & $-0.742$ & $-0.130$ & $\bm{-0.612}$ $[-0.789,-0.404]$ & $5.71\times$ \\
\bench{SAD}    (v2, held-out) & 10 & $-0.165$ & $-0.002$ & $\bm{-0.163}$ $[-0.350,-0.032]$ & $67.6\times$ \\
\midrule
\multicolumn{6}{l}{\footnotesize \emph{Cross-check: pooled probe applied to held-out \bench{SAD} (mean over all $n{=}60$, not paired):}}\\
\bench{SAD}    (pooled probe, held-out)& 60 & $-0.211$ & --- & --- & --- \\
\multicolumn{6}{l}{\footnotesize Closely matches the judge-level held-out shift on \bench{SAD}+\bench{AISI} ($-0.176$, Table~\ref{tab:reverse-canonical}).}\\
\bottomrule
\end{tabular}
\caption{Paired bank-specificity (pooled probe, Qwen3.5-27B layer 12,
CV $0.850$). Both causal-$\Delta$ CIs exclude $0$. The middle separator
shows the same pooled probe applied to the full held-out \bench{SAD}
pool (no matched control on this subset); its activation-level shift of
$-0.211$ converges with the judge-level held-out shift of $-0.176$ from
Table~\ref{tab:reverse-canonical}.}
\label{tab:bank-specificity-paired}
\end{table}

\paragraph{v1 vs v3: the preservation--magnitude trade-off.}
A naive ungated rewriter (v1, no preservation gate) achieves a \emph{larger} apparent
judge shift than the gated v3 pipeline, at the cost of destroying the task.
Table~\ref{tab:v1-v3-tradeoff} quantifies the trade-off on
\bench{AgentHarm}, our hardest case.

\begin{table}[t]
\centering\small
\setlength{\tabcolsep}{3pt}
\begin{tabular}{lrrr}
\toprule
Pipeline & Judge $\Delta$ \Peval & Refusal-class & Activation\\
         &  (\bench{AgentHarm})  & preservation  & probe $\Delta$\\
\midrule
v1 (ungated rewrite) & $\approx -0.55$ & $\sim 10\%$ & $-0.416$ avg.\\
v3 (preservation-gated) & $-0.391$ & $100\%$        & $-0.167$ avg.\\
\midrule
Gap (v1 -- v3) & $\Delta$ judge $-0.16$ & $90\,\text{pp}$ & $\Delta$ probe $-0.25$ \\
\bottomrule
\end{tabular}
\caption{Trade-off between rewrite aggressiveness and task preservation, on
\bench{AgentHarm}. The ungated v1 rewriter produces an apparent
$\Delta \Peval$ that is $\sim 0.16$ larger in magnitude than v3, and a probe
shift that is $\sim 0.25$ larger; but it destroys $90\,\text{pp}$ of
refusal-class preservation. We read this as: $\sim 0.16$/$0.25$ of the v1
``effect'' on safety benchmarks is attributable to inadvertent task drift in
ungated rewriting, not to bank-cue removal. Activation $\Delta$ values are
mean across the 4 Qwen models in Table~\ref{tab:probe-deltas} for per-bench
probes.}
\label{tab:v1-v3-tradeoff}
\end{table}

\paragraph{Cross-benchmark asymmetry: capability vs safety.}
On the \bench{Needham} source pool (capability, donor-restricted; see
Table~\ref{tab:reverse-canonical} row 1, ratio column), v3 reverse with
matched neutral controls gives a bank/control ratio of $7.9\times$
(raw shift $-0.391$ vs control shift $-0.045$). On safety benchmarks
the activation-level ratio is much closer to $1\times$
(Table~\ref{tab:bank-specificity-paired}: \bench{Hawthorne} matched-control shift
$-0.130$ vs bank shift $-0.742$; the ratio there is $5.7\times$ because the
control is not as nearly-null as on capability benchmarks). We interpret
this asymmetry as evidence that on safety benchmarks the eval signal is
substantially \emph{intrinsic} to the test target.

\paragraph{Scaffolding is benchmark-dependent: SAD v3 rejects $60/60$.}
A telling result that strengthens the bimodal reading: when we apply
the v3 preservation gate to the held-out \bench{SAD} pool ($n{=}60$),
\textbf{zero rewrites pass the gate}. Inspecting the gate's rejection
reasons, the pattern is consistent: \bench{SAD} probes situational
awareness via specific MCQ format, second-person ``you'' framing, and
behavior-target labels (e.g.\ ``does the assistant choose to make a
future commitment? Options A/B''). The rewriter, asked to remove
scaffolding, ends up removing the framing that constitutes the test
itself --- and the preservation gate correctly rejects $100\%$ of these.
This is a positive finding for the \emph{gate} (it identifies
non-removable structure) and a negative finding for the
\emph{bank-as-checklist methodology applied to situational-awareness
benchmarks}: on \bench{SAD}, the surface and the test are the same
thing. We read this as evidence that ``scaffolding'' is
benchmark-dependent: removable surface on capability benchmarks IS
the test on situational-awareness benchmarks, generalizing the
capability-vs-safety bimodal into a three-way reading
(capability~$\gg$~safety~$\gg$~situational-awareness in terms of how
much eval-identifiability is surface-removable).

\paragraph{Bank-precision confound on safety benchmarks.}
A reviewer-anticipated concern: our \bench{Hawthorne} bank passes manual
strict-validation at only $14\%$ (Table~\ref{tab:bank-audit}), meaning
$\sim$$86\%$ of \bench{Hawthorne} templates may be content cues rather than
pure scaffolding. The ``safety eval signal is intrinsic'' reading could
therefore be partly confounded with ``our \bench{Hawthorne} bank is
mostly bad.'' We do not fully resolve this confound in the current
submission; the high-precision subset is well-defined
(Appendix~\ref{app:high-precision}) but the v3 reverse pipeline has not
been re-run on the subset because the per-benchmark high-precision count
on \bench{Hawthorne} is $n{=}1$ template, below the statistical floor for
a stable per-category Wilcoxon test. We flag this as the most important
open methodological question for camera-ready, on equal footing with
expanding the manual audit to $n\geq 100$ (currently $n{=}42$). The
robustness of the bimodal capability-vs-safety reading should be read in
light of this caveat.

% ============================================================================
\section{Activation Probes: Within-Domain Separability,
Weak Cross-Domain Transfer}
\label{sec:probes}

\paragraph{Setup.}
We train a linear logistic-regression probe on residual-stream activations
of six open-weight models spanning three model families (Qwen3 8B, 32B;
Qwen3.5 27B, 122B-A10B fp8; \textbf{Mistral-Small-24B-Instruct-2501} and
\textbf{Gemma-2-27B-it} as cross-family checks) to classify orig-eval
vs orig-deploy transcripts at each captured layer, selecting the best
layer per benchmark by 5-fold CV accuracy. We then apply the probe to
v3-rewritten transcripts and to matched controls. \textbf{We treat the
probes as supporting diagnostics, not as evidence of a universal
eval-awareness representation.}

\paragraph{Two metrics, both reported.}
A per-benchmark probe is trained on one benchmark's E vs D and applied to
that benchmark's v3-removed. A pooled probe is trained on the union of
three benchmarks (\bench{Needham}, \bench{Hawthorne}, \bench{MASK})'s E vs D
and applied to each. Both quantify the bank's effect at the activation
level, but at very different magnitudes (Table~\ref{tab:probe-deltas}).

\begin{table}[t]
\centering\small
\setlength{\tabcolsep}{3pt}
\begin{tabular}{lrr}
\toprule
Model & per-bench $\Delta$ & pooled $\Delta$ \\
\midrule
Qwen3-8B               & $-0.103$ & $\bm{-0.535}$ \\
Qwen3-32B              & $-0.173$ & $\bm{-0.598}$ \\
Qwen3.5-27B            & $-0.222$ & $\bm{-0.575}$ \\
Qwen3.5-122B-A10B      & $-0.155$ & $\bm{-0.529}$ \\
\midrule
\textbf{Mistral-Small-24B (cross-family)} & --- & $\bm{-0.289}$ \\
\textbf{Gemma-2-27B-it (cross-family)}    & $-0.269$ & $\bm{-0.250}$ \\
\bottomrule
\end{tabular}
\caption{Bank-removal effect at the activation level, in two probe
constructions. Per-bench probes are trained on each benchmark's own
$E/D$ split; pooled probes are trained on the union of three benchmarks.
Both metrics agree in direction; the pooled-probe magnitude is the
primary number for the supporting-diagnostic claim. We make no monotonic
scaling claim: the 122B-A10B per-bench $\Delta$ is smaller than the
35B-A3B value (Appendix~\ref{app:probes-supplementary}). The two
cross-family rows (Mistral-24B and Gemma-2-27B-it) were added as
robustness checks against the Qwen-family confound: Mistral pooled
CV-acc $0.916$ at layer $12$ (highest of the six models), Gemma pooled
$0.890$ at layer $8$; both show $\sim$$-0.25$ to $-0.29$ pooled $\Delta$
driven by \bench{Hawthorne} ($-0.76$ to $-0.80$), and both replicate the
M1 fail (off-diagonal AUC $0.549$ Mistral, $0.579$ Gemma). This removes
our prior cross-family deferral.}
\label{tab:probe-deltas}
\end{table}

\paragraph{Probes do not transfer cross-benchmark (negative result,
deliberately tested).}
A reviewer-anticipated worry is that our per-benchmark probes are
``benchmark-identity'' detectors rather than eval-vs-deploy detectors.
We tested this by applying each per-benchmark probe to every other
benchmark's E-vs-D classification (Table~\ref{tab:probe-transfer}).
Across all 6 models (four Qwen + cross-family Mistral-Small-24B + Gemma-2-27B-it),
the off-diagonal mean AUC is \textbf{$0.54$--$0.58$ --- essentially chance}.
The M1 fail pattern replicates cleanly cross-family: Mistral off-diag $0.549$,
Gemma off-diag $0.579$ --- statistically indistinguishable from the four Qwen
models at $0.54$--$0.57$. Per-benchmark probes capture benchmark-distribution
structure, not a universal eval-awareness direction.

\paragraph{Shuffled-label null (Qwen3-32B, $n_{\text{perm}}{=}100$).}
A skeptic might still ask: is off-diagonal AUC $\sim$$0.54$ truly at
chance, or modestly above? We computed the null distribution explicitly
on Qwen3-32B by retraining $100$ shuffled-label per-benchmark probes
(stratified label permutation preserves class balance) and evaluating
each on the original target benchmark's E vs D. The actual pooled
off-diagonal AUC is $\bm{0.543}$; the null mean is $\bm{0.495}$; the
actual lands at the $\bm{63\text{rd}}$ percentile of the null
distribution. \textbf{No off-diagonal pair's actual AUC exceeds the
$95$th percentile of its null distribution.} The off-diagonal transfer
result is therefore \emph{statistically indistinguishable from a
shuffled-label baseline} --- a stronger negative result than ``near
chance'' on its face would suggest. The pooled-probe mitigation in the
next paragraph remains the answer.

\begin{table}[t]
\centering\small
\setlength{\tabcolsep}{3pt}
\begin{tabular}{lrr}
\toprule
Model & self mean AUC & off-diag mean AUC \\
\midrule
Qwen3-8B            & 0.989 & $\bm{0.536}$ \\
Qwen3-32B           & 0.989 & $\bm{0.542}$ \\
Qwen3.5-27B         & 0.989 & $\bm{0.567}$ \\
Qwen3.5-122B-A10B   & 0.989 & $\bm{0.538}$ \\
\bottomrule
\end{tabular}
\caption{Cross-benchmark probe transfer (M1). Per-benchmark probes
classify their own benchmark's E-vs-D at AUC $\approx 1.0$ but fail to
transfer (off-diagonal $\approx$ chance). Self AUCs are 5-fold
stratified-CV means at the best layer per benchmark (not held-out test
set); the four-model identical $0.989$ figure reflects this CV pooling.
A held-out-test variant on Qwen3.5-27B at the same layer gives self
AUC $0.95$--$0.99$ on the three benchmarks (Appendix~\ref{app:probes-supplementary}),
indicating the CV figure is not a CV-leakage artifact. Full $3\times 3$
matrices per model are in Appendix~\ref{app:probes-supplementary}.}
\label{tab:probe-transfer}
\end{table}

\paragraph{Mitigation: pooled training partially recovers transfer.}
A pooled probe trained on the union of three benchmarks achieves CV
accuracy $0.83$--$0.85$ and produces v3 $\Delta$ comparable to per-bench
probes (Table~\ref{tab:probe-deltas}, ratio $0.90$--$1.18\times$ across
models). Applied to the never-seen \bench{SAD} pool, the pooled probe
gives a clean $\Delta = -0.211$ that matches the judge-level held-out
$\Delta$ of $-0.176$ from \S\ref{sec:reverse}; per-benchmark probes are
incoherent on \bench{SAD} (one positive, one near-zero, one strongly
negative). Pooled training is the right probe construction for the
within-domain bank-removal claim, but the absolute magnitude remains a
property of the union of these three training distributions, not of a
universal eval-awareness representation.

\paragraph{Cross-DOMAIN transfer (capability vs safety).}
We additionally trained a \emph{capability-only} probe (on \bench{Needham}
$E/D$) and a \emph{safety-only} probe (on \bench{Hawthorne}+\bench{MASK}
$E/D$). Cross-domain transfer is asymmetric and weak: capability$\to$safety
mean AUC $0.543$; safety$\to$capability $0.651$. The probe is
\emph{benchmark-domain-specific}.

\paragraph{Probes are not explained by simple safety-keyword density.}
The training pool is $73\%$ safety / $27\%$ capability by sample count.
If the probe were a safety-keyword detector, this would be the obvious
failure mode. We compute a per-transcript safety-keyword density score
($\sim 40$ keyword patterns: \texttt{refus, harm, danger, ethic, illegal,
weapon, deceive, manipulat, \ldots}) and correlate with the probe's
\Peval{} across $1{,}350$ samples spanning $3$ benchmarks and $3$ probes:
Pearson $r = \bm{+0.009}$, Spearman $\rho = +0.022$. This rules out crude
lexical safety-content confounds. It \emph{does not} rule out semantic
safety/refusal structure that a keyword list would miss; the more rigorous
test for semantic safety confound is the cross-domain transfer above.

\paragraph{Comparison to the stronger probe of
\citet{abdelnabi2025hawthorne}.}
Abdelnabi \& Salem train probes on paired real-vs-hypothetical prompts
(same task, only the URL changes) and label reasoning-trace snippets in
which the model itself reasons about being tested. This design isolates
\emph{model-internal} eval cognition more cleanly than our
whole-transcript input-side probes can. Adopting their paired-prompt design
on our six source benchmarks is the natural next step (\S\ref{sec:limitations}).

% ============================================================================
\section{Robustness Checks}
\label{sec:robustness}

\paragraph{Lexical leave-one-benchmark-out baseline.}
We train a logistic regression on binary word-presence features
(\texttt{CountVectorizer}, $\text{max\_features}\!\in\!\{35, 100, 500, 2000\}$)
to classify E vs D within a benchmark, and re-evaluate under
leave-one-benchmark-out (LOBO). Within-benchmark CV AUC is $0.86$--$0.94$;
LOBO AUC collapses to $0.47$ (chance) across all four feature counts.
Lexical features are highly benchmark-specific. The bank's
cross-benchmark generalization (forward $11/11$; held-out $-0.176$) is
not captured by surface vocabulary.

\paragraph{Cross-extractor bank coverage.}
We extracted a second \bench{Needham} bank using Claude-Sonnet ($96$
templates) and compared to our GPT-extracted subset ($65$ templates).
Both extractors find all six categories. Under a third-party
gemini-2.5-flash semantic-match judge, $86.5\%$ of Claude templates have a
matching GPT template ($26.2\%$ in the reverse direction, dominated by
matching-granularity asymmetry: GPT templates tend to be more fine-grained
and get covered en-masse by broader Claude templates). The high-level
category taxonomy is extractor-agnostic.

\paragraph{Ungated vs gated rewriter (v1 vs v3).}
We ran the same rewriter (\texttt{gemini-2.5-flash}) without our
preservation gate (``v1''). On \bench{AgentHarm}, ungated rewrites
produced $\sim 10\%$ refusal-class preservation; gated rewrites
produced $80.8\%$ (Table~\ref{tab:preservation}) --- an
$\sim 8\times$ improvement. At the activation-probe level,
the ungated and gated effects differ by approximately $0.25$ across four
models. The gap is consistent with inadvertent task drift in ungated
rewrites; the preservation gate isolates the bank-specific component.
\textbf{We do not claim this is PRE-equivalent}: we did not run
\citeauthor{pre2025stealtheval}'s rewriter prompt or dataset; this is
our own ungated baseline.

% ============================================================================
\section{Discussion}
\label{sec:discussion}

\paragraph{Surface vs intrinsic decomposition.}
The forward effect (raw $+0.21$ across $11$ cells;
\textbf{bank-specific causal $+0.192$ pooled across all $11$ cells,
$11/11$ direction-positive, $7/11$ individually significant}) and the
reverse causal effect ($-0.346$ on source after matched control,
$-0.176$ on held-out) together suggest a decomposition: a relatively
small surface component that is removable by preservation-validated
ablation, and a larger intrinsic component entangled with the test
target. The bank-specificity asymmetry across benchmark types
(\S\ref{sec:specificity}) sharpens this: capability evals are
\emph{partly} sanitizable by surface edits; safety evals are not, because
their surface signal and their intrinsic test target are largely the
same thing.

\paragraph{Compositional behavior of cues.}
A natural question is whether cues compose additively. A preliminary
analysis on \bench{Hawthorne} (Appendix~\ref{app:compositionality})
indicated that combining all templates produces a near-zero net
\Peval{} shift while individual solo effects are large. We do
\emph{not} present this as a primary claim --- the combined-injection
input lies outside the probe's training distribution and the per-category
direction effects we observed can plausibly be probe-OOD artifacts.

\paragraph{Implications for benchmark designers.}
Decontamination effort against surface test smells is worth prioritizing
on capability benchmarks (where the bank does targeted work) and is
roughly cosmetic on safety benchmarks (where the eval signal is task-bound).
For safety-eval design specifically, the implication is that surface
sanitization is not a substitute for measuring the model's actual
behavior under deployment-like presentation; methodologies like
\citeauthor{abdelnabi2025hawthorne}'s paired-prompt design are required.

% ============================================================================
\section{Conclusion}
\label{sec:conclusion}

We introduce \emph{test smells}: a 258-template bank of canonical evaluation cues
extracted from six widely-used LLM benchmarks. Three convergent measurements ---
forward injection ($+0.21$ raw across $11/11$ cells; bank-specific causal
$+0.192$ across all $11/11$ control-scored cells, $11/11$ direction-positive,
$7/11$ individually significant),
preservation-gated reverse intervention ($-0.346$ causal effect on source
/ $-0.176$ on never-seen \bench{SAD}+\bench{AISI}), and convergent
preservation oracles at $73$--$100\%$ across capability and safety benchmarks
--- support a decompositional reading: eval-identifiability has a
\emph{removable surface component} (the bank, especially on capability benchmarks
where it dominates matched neutral rewrites $7.9\times$) and a \emph{larger
intrinsic component} entangled with the test target, particularly on safety
and situational-awareness benchmarks (on \bench{SAD}, $60/60$ v3 rewrites are
rejected by the gate: surface and test are the same thing). Activation probes on six open-weight models spanning three families (Qwen3/Qwen3.5, Mistral, Gemma-2) corroborate the
direction at the representation level but only under cross-benchmark pooled
training; per-benchmark probes fail to transfer (off-diagonal AUC $\sim 0.54$,
shuffled-label-null indistinguishable),
and we explicitly do not claim a universal eval-awareness representation. A
$n{=}60$ dual-annotator safety spot-check ($\kappa{=}0.509$) finds $\bm{0/30}$
v3 task unlocks under either annotator, corroborating the preservation oracles. Our
recommendation to benchmark designers: surface decontamination is worth
prioritizing on capability benchmarks and is roughly cosmetic on safety
benchmarks, where the eval signal is task-bound by design.

% ============================================================================
\section{Limitations}
\label{sec:limitations}

We organize the principal limitations as six named items. Each names
the limitation, summarizes what we did to mitigate it in this submission,
and points at the camera-ready follow-up.

\paragraph{L1. Single primary judge.}
All headline causal effects flow through \texttt{gpt-5-mini}. We
mitigate with two formal checks (\S\ref{sec:preservation}): a dual-judge
agreement study on $n{=}153$ stratified samples
($\bm{\kappa = 0.413}$, $92.2\%$ agreement,
Table~\ref{tab:dual-judge-perbench}) and a position-bias 2-swap on
$n{=}81$ samples ($\bm{6.2\%}$ flip rate). Both indicate the judge
is moderately stable beyond chance. Neither fully eliminates
single-judge dependence; a full causal re-run under a different
primary judge is the camera-ready follow-up.

\paragraph{L2. Manual bank validation is $n{=}42$, with \bench{Hawthorne}
at $14\%$ strict pass.}
Aggregate manual validation (Table~\ref{tab:bank-audit}) is
$67\%$ strict / $88\%$ inclusive at $n{=}42$
($n{=}7$ per benchmark). \bench{Hawthorne} validates at only
$14\%$ ($1/7$), implying the \bench{Hawthorne} bank carries
content cues alongside scaffolding. The bimodal capability-vs-safety
reading should be interpreted in light of this. Mitigation in this
submission: the high-precision subset is defined in
Appendix~\ref{app:high-precision} and the gap is flagged explicitly in
\S\ref{sec:specificity}. Camera-ready: expand audit to $n\geq 100$ and
re-run v3 on the high-precision subset.

\paragraph{L3. Per-cell forward bank-specificity is noisy.}
The pooled forward causal $\Delta = +0.192$ across $11/11$ cells
($11/11$ direction-positive) is supported at the aggregate level, but
only $7/11$ cells are individually significant. The $4$ borderline cells
are exactly those with small per-cell $n$ ($15$--$23$). Camera-ready:
scale per-cell $n$ to $50$+ for tighter individual CIs.

\paragraph{L4. SAD v3 gate rejects $60/60$.}
Running v3 on the held-out \bench{SAD} pool rejected all $60$ rewrites
(\S\ref{sec:specificity}). We treat this as a positive finding
(``scaffolding is benchmark-dependent'') and the held-out judge result
($-0.176$, $p<10^{-13}$, Table~\ref{tab:reverse-canonical}) therefore
uses v2 not v3 on the held-out pool. Camera-ready: construct a
SAD-specific rewriter that preserves the MCQ-format scaffolding while
removing only narrowly-targeted cues.

\paragraph{L5. Probes are benchmark-domain-specific.}
Activation probes (\S\ref{sec:probes}) are explicitly supporting
diagnostics, not evidence for a universal eval-awareness representation.
Per-benchmark probes do not transfer cross-benchmark (off-diagonal
AUC $\sim 0.54$--$0.58$ across all six models, including both
cross-family probes Mistral-24B at $0.549$ and Gemma-2-27B-it at
$0.579$). A shuffled-label null on Qwen3-32B ($n_{\text{perm}}{=}100$)
confirms this off-diagonal AUC is statistically indistinguishable from
chance (actual $0.543$ vs null mean $0.495$, $63$rd percentile). A
pooled probe rescues the bank-removal effect but is
benchmark-distribution-specific. The stronger probe construction is
\citeauthor{abdelnabi2025hawthorne}'s paired-prompt reasoning-trace
methodology, which we do not implement here. Camera-ready: adopt their
paired-prompt design across our six source benchmarks.

\paragraph{L6. v3 unlock rate on \bench{AgentHarm}: $7.7\%$
(single-judge); $0/30$ in dual-annotator validation.}
The preservation-gated v3 rewriter unlocks $2/26$ \bench{AgentHarm}
samples from \textsc{Safe} to \textsc{Harmful} under the strict
single-judge safety-rubric (Table~\ref{tab:preservation}). On aggregate
v3 is slightly safer ($19.2\%$ harden rate). A stratified $n{=}60$
dual-annotator spot-check (gemini-2.5-flash + claude-haiku-4.5,
$\kappa{=}0.509$, \S\ref{sec:preservation}) finds \textbf{$0/30$ task
unlocks under \emph{either} annotator} on paired
(\bench{AgentHarm}, \bench{Hawthorne}) samples --- the rubric-judge
$7.7\%$ unlock rate may reflect single-judge labeling noise on
borderline samples. We surface both numbers honestly. Camera-ready:
larger dual-annotator $n$ on a fully stratified design, plus a
stricter gate variant that filters on preservation \emph{plus} per-task
safety-rubric agreement.

\paragraph{Compute and reproducibility (not a limitation; included here for completeness).}
All judge / rewriter inference is API-based (gpt-5-mini, gemini-2.5-flash,
claude-haiku-4.5). Activation extraction ran on RunPod H200 and B200
community-cloud pods across six models, accumulating $\sim$$54$
H200-equivalent hours (Qwen3-8B/32B/27B fp16, Qwen3.5-122B-A10B fp8,
plus Mistral-Small-24B and Gemma-2-27B-it both bf16 on a B200) plus
pilot runs. Total inference cost: $\sim$$\$80$ OpenRouter
$+\sim$$\$220$ RunPod $\approx\$300$ total. We release the bank,
paired datasets, and analysis code at
\href{https://anonymous.4open.science/}{anonymous.4open.science/[anonymized]}.

% ============================================================================
\section{Ethics}
\label{sec:ethics}

\paragraph{Dual use.}
Documenting concrete test smells could in principle help an adversarial
actor train or prompt a model to better ``hide'' eval-aware behavior. We
judge the net effect of disclosure to be positive: benchmark designers
need the same information to detect and remove cues, and obscurity over
the test-smell taxonomy provides no robust defense (the cues we describe
are already widely visible in published benchmarks). This mirrors prior
norms in red-teaming and prompt-attack disclosure
\citep{perez2023redteaming}.

\paragraph{Data provenance.}
Source transcripts inherit their upstream benchmarks' licenses. The
\bench{Needham} pool is gated and accessed under research-use terms; we
do not redistribute it. The naturalistic chat in our recipient pools
derives from \bench{WildChat} \citep{zhao2024wildchat}. Our released
artifacts are distributed under CC-BY-NC-SA-4.0, matching the most
restrictive upstream license, with a dataset card following
\citet{gebru2021datasheets} and \citet{pushkarna2022datacards}.

\paragraph{AI assistance disclosure.}
Per ARR policy, GPT-5.1 and GitHub Copilot were used for code authoring,
debugging, and prose proofreading. Substantive content, claims, analysis
design, and final writing were authored by the human authors.

% ============================================================================
\bibliographystyle{acl_natbib}
\bibliography{refs}

% ============================================================================
\newpage
\appendix

\section{Bank: per-benchmark $\times$ per-category counts}
\label{app:bank}

\begin{table}[h]
\centering\small
\setlength{\tabcolsep}{3pt}
\begin{tabular}{lrrrrrrr}
\toprule
Benchmark & SI & TF & AF & FS & CT & EP & Total \\
\midrule
\bench{Needham}    & 19 & 10 & 9  & 4  & 17 & 6  & 65 \\
\bench{Hawthorne}  & 0  & 0  & 0  & 4  & 4  & 7  & 15 \\
\bench{MASK}       & 6  & 1  & 1  & 6  & 3  & 10 & 27 \\
\bench{AgentHarm}  & 2  & 2  & 4  & 6  & 5  & 7  & 26 \\
\bench{BBH}        & 2  & 8  & 2  & 10 & 23 & 0  & 45 \\
\bench{MALT}       & 15 & 21 & 15 & 5  & 23 & 1  & 80 \\
\midrule
\textbf{Total}     & \textbf{44} & \textbf{42} & \textbf{31} & \textbf{35} & \textbf{75} & \textbf{31} & \textbf{258} \\
\bottomrule
\end{tabular}
\caption{Bank composition by source benchmark and category.
SI=\category{SystemInstructions}, TF=\category{TaskFormatting},
AF=\category{AnswerFormatting}, FS=\category{FictionalScenario},
CT=\category{ContrivedTask}, EP=\category{EthicalProbe}.
\bench{MALT} contributes the largest share ($80$ templates, $31\%$);
\bench{Hawthorne} the smallest ($15$, $6\%$). The \bench{Hawthorne} bank
notably contains zero \category{SystemInstructions} or \category{TaskFormatting}
templates --- consistent with safety harmful-request benchmarks lacking the
scaffolding tokens that capability/MCQ benchmarks expose.}
\label{tab:bank-by-benchmark}
\end{table}

\section{High-precision bank subset}
\label{app:high-precision}

We define the \emph{high-precision subset} as those templates whose
per-benchmark stratified manual-validation pass rate is $\geq 80\%$
strict (Table~\ref{tab:bank-audit}). Under this rule, the high-precision
subset is:

\begin{itemize}\itemsep0pt
\item \bench{BBH}: all $45$ templates ($100\%$ strict pass)
\item \bench{Needham}: all $65$ templates ($100\%$ strict pass)
\item \bench{MALT}: $\sim 80$ templates ($86$--$100\%$ inclusive pass)
\item \bench{AgentHarm}: bottom-half of templates excluded ($\sim 15$ retained
  out of $26$, $57\%$ pass rate; subset definition requires per-template audit
  beyond stratified $n{=}7$)
\item \bench{MASK}: $\sim 12$ retained out of $27$ ($43\%$)
\item \bench{Hawthorne}: $\sim 2$ retained out of $15$ ($14\%$) --- the
  high-precision \bench{Hawthorne} subset is too small for a stable
  per-category Wilcoxon test in the v3 reverse pipeline. We flag this as
  the most important limitation of our bimodal capability-vs-safety reading
  (\S\ref{sec:specificity}, \S\ref{sec:limitations}).
\end{itemize}

The full main effects (\S\S\ref{sec:forward}--\ref{sec:reverse}) are
reported on the full bank. Re-running on the high-precision subset is
straightforward via \texttt{scripts/high\_precision\_subset.py} in our
released code and is the natural camera-ready follow-up.

\section{Reverse intervention: per-category breakdown}
\label{app:reverse-cat}

\begin{table}[h]
\centering\small
\begin{tabular}{lrrrrrr}
\toprule
Category & $n$ & base & treat & ctrl & raw $\Delta$ & causal $\Delta$ \\
\midrule
\category{ContrivedTask}        & 57 & 0.922 & 0.494 & 0.866 & $-0.428$ & $\bm{-0.373}$ \\
\category{TaskFormatting}       & 43 & 0.909 & 0.473 & 0.838 & $-0.436$ & $\bm{-0.366}$ \\
\category{SystemInstructions}   & 28 & 0.906 & 0.568 & 0.880 & $-0.338$ & $\bm{-0.312}$ \\
\category{EthicalProbe}         & 19 & 0.853 & 0.530 & 0.870 & $-0.323$ & $\bm{-0.340}$ \\
\category{AnswerFormatting}     & 11 & 0.973 & 0.652 & 0.926 & $-0.321$ & $\bm{-0.274}$ \\
\category{FictionalScenario}    &  4 & 0.909 & 0.684 & 0.958 & $-0.225$ & $\bm{-0.275}$ \\
\midrule
\textbf{All}                    & \textbf{93} & 0.900 & 0.509 & 0.855 & $\bm{-0.391}$ & $\bm{-0.346}$ \\
\bottomrule
\end{tabular}
\caption{v3 reverse intervention with matched neutral controls, per
bank category. Wilcoxon $p<10^{-4}$ for the all-categories comparison
and for the four categories with $n\geq 19$.}
\label{tab:reverse-by-cat}
\end{table}

\section{Reverse intervention: sub-benchmark composition of the
\bench{Needham} source pool}
\label{app:reverse-bench}

The $n{=}93$ paired samples in Table~\ref{tab:reverse-canonical} row 1 come
from the \bench{Needham} donor pool, which is itself an aggregate of
$\sim 20$ sub-benchmarks. The dominant sub-benchmarks (by count in the
$n{=}161$ pre-eligibility pool, of which $n{=}93$ pass the $P_b<0.5$
flip-eligibility filter) are:

\begin{table}[h]
\centering\small
\begin{tabular}{lr}
\toprule
Sub-benchmark & $n$ in v3 reverse pool \\
\midrule
aider (coding-eval)            & 27 \\
cybench (CTF-eval)             & 25 \\
agentharm (safety-tool)        & 22 \\
smol (agent-eval)              & 14 \\
gandalf, mosscap (jailbreak)   & 22 \\
wmdp (knowledge-eval)          & 9 \\
gaia, mmlu (general)           & 12 \\
mvp, swe(bench), model         & 20 \\
ethics, llm, aya, kobbq, moralexceptqa, safetybench & 8 \\
\midrule
\textbf{Total (pre-eligibility)} & \textbf{161} \\
\textbf{Total (post-eligibility, $P_b{\geq}0.5$)} & \textbf{93} \\
\bottomrule
\end{tabular}
\caption{Sub-benchmark composition of the \bench{Needham} source pool.
\bench{Needham} aggregates many capability and safety sub-benchmarks; the
$7.9\times$ bank/control ratio in Table~\ref{tab:reverse-canonical} is on
the pooled mixture.}
\label{tab:reverse-subbench}
\end{table}

\section{Preservation oracles: per-task BBH breakdown}
\label{app:bbh-pertask}

\begin{table}[h]
\centering\small
\setlength{\tabcolsep}{3pt}
\begin{tabular}{lrrrr}
\toprule
BBH task & $n$ orig & orig acc & v3 acc & $\Delta$ \\
\midrule
boolean\_expressions & 19 & 1.00 & 1.00 & 0 \\
multistep\_arithmetic & 19 & 1.00 & 1.00 & 0 \\
object\_counting & 18 & 1.00 & 1.00 & 0 \\
date\_understanding & 1 & 1.00 & 1.00 & 0 \\
hyperbaton & 14 & 0.86 & 0.93 & $+0.07$ \\
logical\_deduction (3-obj) & 7 & 1.00 & 1.00 & 0 \\
logical\_deduction (5-obj) & 10 & 1.00 & 0.30 & $-0.70$ \\
logical\_deduction (7-obj) & 11 & 0.73 & 0.18 & $-0.55$ \\
disambiguation\_qa & 6 & 0.50 & 0.50 & 0 \\
ruin\_names & 8 & 0.88 & 0.63 & $-0.25$ \\
sports\_understanding & 20 & 0.75 & 0.50 & $-0.25$ \\
\dots (13 other tasks, mostly preserving) & --- & --- & --- & --- \\
\midrule
\textbf{Aggregate (24 tasks)} & 252 & 0.857 & 0.770 & $-0.087$ \\
\bottomrule
\end{tabular}
\caption{BBH per-task accuracy preservation. The aggregate $89.8\%$
relative preservation hides per-task variation: most tasks preserve at
$100\%$, but $\bench{logical\_deduction}$ (5- and 7-object variants) loses
substantial accuracy under v3, suggesting that for those tasks the
scaffolding tokens (e.g.\ enumeration formatting) are partly load-bearing
for the model's answer-extraction.}
\label{tab:bbh-pertask}
\end{table}

\section{Activation probes: supplementary}
\label{app:probes-supplementary}

The full $3\times 3$ M1 cross-benchmark probe transfer AUC matrices per
model are included in the release at
\texttt{results/activation\_probing/m1\_cross\_transfer\_\_\\\$model.json}.
A representative model (Qwen3.5-27B):

\begin{table}[h]
\centering\small
\setlength{\tabcolsep}{3pt}
\begin{tabular}{l|rrr}
\toprule
Source probe $\backslash$ Target & \bench{Needham} & \bench{Hawthorne} & \bench{MASK} \\
\midrule
\bench{Needham}    & 1.00 & 0.45 & 0.59 \\
\bench{Hawthorne}  & 0.53 & 1.00 & 0.55 \\
\bench{MASK}       & 0.60 & 0.55 & 1.00 \\
\bottomrule
\end{tabular}
\caption{Cross-benchmark probe transfer AUC, Qwen3.5-27B. Off-diagonal
mean $0.54$, essentially chance. Per-bench probes are benchmark-identity
detectors, not eval-vs-deploy detectors.}
\label{tab:m1-27b}
\end{table}

Pooled-probe per-layer CV curves for the same model peak at layer $12$
($\text{CV}=0.850$); we use this for the bank-specificity computations
in \S\ref{sec:specificity}. Per-model best-layer depth-fraction
exploratory data (layer-salience) is in the released results JSON; we
do not foreground it as a paper claim because our whole-transcript probes
are not directly comparable to \citet{xue2025linearcontroleval}'s
paired-prompt reasoning-trace setup.

\section{Forward neutral-injection: per-sample data}
\label{app:fwd-neutral-supplementary}

The per-sample triples $(P_{\text{ctrl}}, P_{\text{bank}}, P_{\text{neutral}})$
and bootstrap distributions are shipped at
\texttt{results/forward\_neutral\_injection.json} ($n{=}20$ per benchmark on
\bench{Hawthorne} and \bench{MASK}). The headline summary is in
Table~\ref{tab:fwd-neutral}; per-sample data supports the bootstrap CIs reported
there ($B{=}1000$). Scaling to $n{=}100$ per benchmark is the natural
camera-ready follow-up; on each of the other $9$ forward injection cells
($+0.21$ raw mean, Table from~\S\ref{sec:forward}) we have only the raw
shift, not the neutral-control-adjusted bank-specific component.

\section{Lexical baseline: sweep over feature counts}
\label{app:lexical}

\begin{table}[h]
\centering\small
\begin{tabular}{lrr}
\toprule
$\text{max\_features}$ & Within-bench AUC & LOBO AUC \\
\midrule
35     & 0.82 & 0.44 \\
100    & 0.86 & 0.46 \\
500    & 0.91 & 0.47 \\
2000   & 0.94 & 0.47 \\
\bottomrule
\end{tabular}
\caption{Lexical leave-one-benchmark-out baseline. Within-benchmark
$5$-fold CV AUC vs LOBO AUC. Adding features helps within a benchmark
($+0.12$ AUC); it does not help across benchmarks ($+0.03$).}
\label{tab:lexical-lobo}
\end{table}

\section{Cross-extractor bank category mapping}
\label{app:cross-extractor}

We re-extracted the \bench{Needham} subset of the bank with a different
LLM extractor (Claude-Sonnet) and compared to our GPT-extracted
\bench{Needham} subset. Both extractors recover all six categories.
The aligned-template counts per category (per
\texttt{results/pri5\_bank\_overlap.json}):

\begin{itemize}\itemsep0pt
\item \category{SystemInstructions}: $5$/$5$ Claude templates align with GPT
\item \category{TaskFormatting}: $4$/$5$ align
\item \category{AnswerFormatting}: $3$/$4$ align
\item \category{FictionalScenario}: $2$/$3$ align
\item \category{ContrivedTask}: $7$/$8$ align
\item \category{EthicalProbe}: $2$/$3$ align
\end{itemize}

Aggregate alignment rate: $23/28 = 86.5\%$. Reverse direction (GPT
templates having a matching Claude template) is lower at $26.2\%$, but this
is dominated by matching-granularity asymmetry: GPT-extracted templates
tend to be finer-grained and get covered en-masse by broader Claude
templates. The high-level category taxonomy is extractor-agnostic.

\section{Compositionality (exploratory; not a paper claim)}
\label{app:compositionality}

We attempted a per-template solo vs combined-injection compositionality
analysis on \bench{Hawthorne} (combined injection of all $13$ Hawthorne
templates simultaneously vs each individual injection). The combined-shift
was near-zero ($\Delta \approx +0.04$ on Qwen3.5-27B) while the sum of
solo shifts was large ($\approx -1.43$ on the same model), suggesting
\emph{non-additive} composition. We do \emph{not} promote this to a
paper claim because the combined-injection input lies outside the
per-benchmark probe's training distribution (M1:
off-diagonal AUC $\sim 0.54$), and the per-category direction effects
we observed plausibly reflect probe out-of-distribution behavior rather
than a genuine non-additivity. The analysis is included for completeness.

\section{Compute and cost}
\label{app:compute}

Total inference: $\sim$$\$80$ across activation extraction, judge calls,
behavioral measurement, lexical baselines, and cross-extractor analysis.
H200 GPU time $\sim$$10$ hours. OpenRouter API calls $\sim$$1{,}500$ for
judge + rewriter + neutral-prefix + safety-class judges.

\end{document}
